\chapter{Kết luận}

\section{Tổng kết}

Khóa luận này xuất phát từ một quan sát đơn giản về cách các hệ thống điều phối xe đi chung
đang vận hành: chúng báo cho hành khách một giờ hẹn, rồi tiếp tục tối ưu lại và dời giờ ấy nhiều
lần, trong khi không một thước đo chuẩn nào của ngành ghi nhận điều đó. Câu hỏi đặt ra là liệu
một hệ thống có thể đưa ra lời hứa mà nó thực sự giữ được hay không, và cái đó đắt bao nhiêu.

Để trả lời, chúng tôi xây dựng cơ chế lời hứa có biên được chứng nhận, với ba thành phần. Thành
phần thứ nhất là phép phân rã ba mốc, dựng một mốc phản thực bằng cách chiếu lại chính kế hoạch
cũ từ vị trí thật hiện tại của xe, nhờ đó tách được phần dịch chuyển do nhà điều hành gây ra
khỏi phần do thế giới vật lý gây ra. Thành phần thứ hai là sổ cái cam kết riêng cho từng hành
khách, cộng dồn theo thời gian và đối chiếu với một trần, được kiểm tra ngay tại thời điểm ra
quyết định. Thành phần thứ ba là hệ thống chứng cứ niêm phong, bảo đảm rằng mọi so sánh giữa
các cấu hình đều kiểm chứng được tới từng lần chạy.

Về mặt lý thuyết, đóng góp trung tâm là việc phân biệt hai cơ sở tính sổ cái và ba định lý đi
kèm. Tính theo hành vi của nhà điều hành thì luôn thỏa được, vì phương án không làm gì theo định
nghĩa không gây ra gì. Tính theo kết quả mà hành khách nhìn thấy thì có thể khiến hệ thống không
còn bất kỳ quyết định hợp lệ nào, vì ngay cả việc không làm gì cũng bị tính phí trôi dạt. Và
quan trọng nhất, cơ sở kết quả không bao giờ phục vụ được nhiều hơn cơ sở hành vi với bất kỳ
trần nào --- đây là một định lý chứ không phải một quan sát may mắn ở một giá trị cụ thể. Phân
biệt này trùng với phân biệt cổ điển giữa nghĩa vụ phương tiện và nghĩa vụ kết quả trong luật
hợp đồng, và điểm mới nằm ở chỗ mốc tham chiếu ngoại sinh \emph{tính được} ranh giới ``ngoài
tầm kiểm soát'' mà luật xử lý bằng văn xuôi.

Về mặt thực nghiệm, 558 lần chạy trên một panel 20 ô đã niêm phong cho bốn nhóm kết quả. Nhóm
thứ nhất là quy trách: $89{,}06\%$ lượng dịch chuyển mà hành khách nhìn thấy trong chế độ không
cam kết là do bộ điều phối tự gây ra, và thiệt hại dồn vào một cái đuôi mỏng --- $23{,}2\%$ số
hành khách hứng toàn bộ, người chịu thiệt nặng nhất bị dời gần mười bốn phút mà hệ thống vẫn
tính là phục vụ thành công. Nhóm thứ hai là đường cong giá của cam kết trên chín cấu hình, với
phát hiện rằng phần đắt nhất là quyết định cam kết chút nào đó chứ không phải độ chặt của lời
hứa về thời gian. Nhóm thứ ba xác nhận cả ba định lý, kể cả dưới dữ liệu tắc đường đo thật của
Manhattan, nơi cơ sở kết quả tê liệt trên 19 trên 20 ô còn cơ sở hành vi không hỏng một ô nào.
Nhóm thứ tư kiểm tra độ vững: chín thước đo cho cùng một thứ hạng, không một ngoại lệ.

Ngoài ra, khóa luận ghi nhận bốn kết quả phụ mà chúng tôi cho là đáng chú ý. Thứ nhất, sáu cơ
chế bảo vệ giờ đón khác nhau --- kể cả cơ chế bảo đảm một chiều mà văn liệu sử dụng --- rơi vào
đúng hai hành vi, vì hình thức của một lời bảo đảm chỉ có ý nghĩa trong chừng mực bộ sinh
phương án có thể thực thi nó. Thứ hai, giá của một lời hứa không phải hằng số mà tăng theo mức
tự do của nhà điều hành, và mặt cong hai chiều ở mục \ref{muc:matcong} cho thấy hai yếu tố này
\emph{tương tác}: trần càng chặt thì việc cho thêm tự do càng đắt, từ $0{,}42$ lên $1{,}44$
điểm phần trăm khi trần siết từ 300 xuống 30 giây. Thứ ba, hệ quả của định lý về hai cơ sở
biểu hiện theo kiểu phân loại chứ không theo kiểu biên --- cơ sở kết quả mất nguyên cả ô chứ
không phục vụ ít hơn vài hành khách --- và lý do là số học: khả thi cấp ô sụp theo $(1-p)^{n}$
với $n$ là số hành khách mỗi ô.

Thứ tư, ba hướng cải tiến bộ sinh phương án mà chúng tôi dự định thực hiện đều bị chính dữ liệu
của hệ thống bác bỏ, mỗi hướng vì một lý do khác nhau, và điều đó tự nó là một kết quả. Xếp hạng
vị trí chèn và cắt tỉa sớm là phép toán rỗng vì $96{,}0\%$ quyết định cấp xe chỉ có đúng một ứng
viên khả thi trên $217.736$ quan sát. Đặt trước ngân sách chạy được nhưng trùng khớp tới từng
quyết định với đối chứng, vì tầng cảnh báo bị đặt bên dưới một tầng mục tiêu liên tục đo cùng
đại lượng ở thang mịn hơn --- ý tưởng ấy đã nằm sẵn trong hệ thống ở dạng mạnh hơn. Và bộ sinh
biết sắp xếp lại, vốn cũng đã có sẵn, va vào một bất biến thứ hai của mô hình về thời điểm một
chặng hết sửa được. Điểm chung của cả ba là chúng được loại bằng phép đo rẻ tiền trên bằng chứng
đã có, chứ không bằng nhiều ngày chạy mô phỏng --- và bằng chứng ấy do chính cơ chế lời hứa sinh
ra.

\section{Hạn chế}

Chúng tôi muốn nêu các hạn chế một cách rõ ràng, vì một số trong đó ảnh hưởng trực tiếp tới
cách đọc kết quả.

Hạn chế lớn nhất là \textbf{quy mô}. Panel 108 yêu cầu và 8 xe mỗi ô nhỏ hơn nhiều so với các
thực nghiệm quy mô Manhattan trong văn liệu. Lựa chọn này có chủ ý --- mục tiêu là so sánh có
kiểm soát giữa các cấu hình trên cùng đầu vào --- nhưng hệ quả là các con số tuyệt đối không so
trực tiếp được với văn liệu.

Hạn chế thứ hai là \textbf{phần trùng lặp với công trình gần nhất}. Như mục \ref{muc:doichieu} đã trình bày chi
tiết, sáu trong số các hạng mục của khóa luận đã có trong Fielbaum và Alonso-Mora (2020): khái
niệm, đo độ lớn, đếm số lần, phân loại nguyên nhân định tính, đường cong đánh đổi, và cơ chế
ratchet. Chúng tôi coi phần đó là nền chứ không phải đóng góp, và đã cố gắng ghi rõ ranh giới ở
bảng \ref{bang:trung}.

Hạn chế thứ ba là \textbf{trôi dạt ở panel chính chỉ là cận dưới}. Bộ dữ liệu thử gốc có đồ thị
thời gian di chuyển tĩnh, nên trôi dạt đo được đến từ độ trễ thực thi chứ không từ tắc đường.
Thí nghiệm bơm trôi dạt đã bù một phần nhưng chưa trọn vẹn: nó cho kết quả về tính khả thi nhưng
không cho kết quả về tỉ lệ phục vụ, vì một bất biến miền của hệ thống chặn lại.

Hạn chế thứ tư là \textbf{lập luận phúc lợi kém chắc chắn hơn phần kỹ thuật}. Nó dựa trên hệ số
định giá từ một bối cảnh địa lý khác và trên giả định rằng mọi chuyến bị mất có cùng thặng dư.
Chúng tôi trình bày nó như một ngưỡng chứ không phải một kết luận.

Cuối cùng, khóa luận \textbf{không tuyên bố vượt trội về tỉ lệ phục vụ}. Mọi cấu hình có cam
kết đều phục vụ ít hành khách hơn cấu hình không cam kết. Đóng góp nằm ở việc biến một đại
lượng trước nay chỉ đo được sau sự việc thành một ràng buộc cưỡng chế được trước khi quyết
định, và ở việc chỉ ra rằng cách tính sổ cái quyết định hệ thống có vận hành nổi hay không.

\section{Hướng phát triển}

Từ những hạn chế trên, chúng tôi đề xuất năm hướng tiếp theo, xếp theo thứ tự ưu tiên mà chúng
tôi cho là hợp lý.

\textbf{Đo Định lý 3.3 trên một panel có ít hành khách hơn mỗi ô.} Mục \ref{muc:troidat} cho thấy phép so
sánh tỉ lệ phục vụ giữa hai cơ sở không lấy được trên panel hiện tại, vì khả thi cấp ô sụp theo
$(1-p)^{n}$ với $n \approx 80$. Lối thoát tự nhiên là giảm $n$: một panel nhiều ô nhỏ hơn, mỗi
ô vài chục hành khách, sẽ giữ được nhiều ô khả thi ở cùng mức trôi dạt, nhờ đó có tập ô chung
để so. Đây là thay đổi về thiết kế panel chứ không phải về hệ thống, nên chi phí thấp.

\textbf{Mở rộng mặt cong sang chiều tự do thứ ba.} Mục \ref{muc:matcong} đã dựng mặt cong trên hai mức tự
do và cho thấy tự do với độ chặt \emph{tương tác} chứ không cộng độc lập. Chiều tự do ở đây là
cửa sổ đón; hai chiều tự nhiên tiếp theo là cỡ đội xe và mật độ nhu cầu. Nếu quy luật giữ nguyên
trên cả ba chiều thì có thể rút ra một quy tắc chọn trần theo mức tự do, thay vì phải đo lại cho
từng cấu hình vận hành.

\textbf{Cửa sổ đón thương lượng có chọn lọc.} Như mục \ref{muc:dathu} đã nêu, hướng này chưa bị bác bỏ. Bài
toán đã được đặt đúng với cả trần lẫn kiểu hỏng đều đã đo: lấy được bao nhiêu trong 172 hành
khách mà vẫn giữ chờ trung vị quanh 407 giây. Tuy nhiên nó đòi hỏi sửa mã nguồn và mở rộng phát
biểu vấn đề sang kiểm soát tiếp nhận, nên cần cân nhắc kỹ về phạm vi.

\textbf{Nới lỏng bất biến miền về cửa sổ đón.} Phát hiện ở mục \ref{muc:batbien} --- rằng hệ thống chưa biểu
diễn được trạng thái mà riêng giao thông phá vỡ một cửa sổ đã cam kết --- là một giới hạn thật
của mô hình. Cơ chế nới lỏng theo mốc ngoại sinh đã tồn tại trong hệ thống cho bộ sinh phương án;
mở rộng nó sang tầng kiểm tra trạng thái sẽ cho phép nghiên cứu chế độ cao điểm một cách đầy đủ.

\textbf{Đóng băng chặng đang chạy, rồi mới mở khóa bộ sinh biết sắp xếp lại.} Mục
\ref{muc:bosinh} cho thấy hệ thống đã có sẵn cơ chế sắp xếp lại nhưng nó hỏng trên 18 trên 20
ô, vì mô hình miền chỉ đóng băng một điểm dừng khi xe \emph{tới nơi}, còn bộ mô phỏng khóa chặng
ngay khi xe \emph{lăn bánh}. Việc cần làm có phạm vi rõ ràng và phải theo đúng thứ tự: thêm
chuyển trạng thái đóng băng chặng đang bay ở tầng miền --- kiểu dữ liệu lộ trình đã cho phép
trạng thái ``đã đóng băng nhưng chưa thực hiện'', chỉ thiếu điểm gọi; cho chính sách vector cứng
nhận hạn mức sửa chữa; và nới điều kiện kiểm tra cấu hình để hai thứ đi được cùng nhau. Chỉ sau
ba bước đó mới đặt được câu hỏi mà mục 4 thật sự muốn hỏi: liệu quyền sắp xếp lại có thu hồi
được những hành khách mà cổng cam kết đang từ chối hay không. Chúng tôi lưu ý rằng ô duy nhất
chạy trọn vẹn cho kết quả bằng đúng nhánh không cam kết, nên kỳ vọng nên đặt thấp.

\section{Kết luận cuối}

Đóng góp mà chúng tôi tin là bền nhất của khóa luận không phải một con số cụ thể, mà là một cách
đặt vấn đề. Văn liệu hiện hành \emph{mô tả} hiện tượng lời hứa bị xê dịch; khóa luận này
\emph{chặn} nó. Sự khác biệt giữa hai việc ấy không phải là mức độ mà là bản chất: thước đo của
văn liệu tham chiếu tới giá trị thực tế, chỉ biết được sau khi hành khách đã được đón, nên về
mặt toán học không thể trở thành ràng buộc; còn cơ chế ở đây định nghĩa trên các đại lượng đều
biết được tại thời điểm quyết định.

Và khi biến một thước đo thành một ràng buộc, một câu hỏi mới xuất hiện mà trước đó không có lý
do gì để hỏi: \emph{tính sổ cái theo cái gì}. Câu trả lời hóa ra quyết định hệ thống có vận hành
nổi hay không --- điều mà ba định lý và thí nghiệm tắc đường thật đã chỉ ra. Cách tự nhiên nhất
để biến thước đo của văn liệu thành ràng buộc chính là cơ sở kết quả, và đó cũng chính là cách
hỏng.
